Single-photon LiDAR can support long-range optical sensing under severe photon limitations, but uniform illumination may waste transmitted energy on low-value regions. This short communication presents a physics-guided adaptive photon-allocation strategy for defense-oriented single-photon LiDAR. The method combines posterior target uncertainty, predicted signal-to-background ratio, measurement history, and exploration to determine the next sensing region. A physically grounded simulator incorporates pulse and dwell accounting, beam divergence, target-footprint coupling, two-way atmospheric attenuation, distributed backscatter, timing jitter, detector efficiency, dark counts, dead time, pile-up, and afterpulsing. The proposed information-guided strategy is compared with uniform, random, greedy-posterior, and entropy-based allocation. Across 30 simulated episodes per method, the proposed strategy achieved a 93.3% detection rate, compared with 90.0% for greedy and entropy allocation, 76.7% for uniform allocation, and 63.3% for random allocation. It also reduced average transmitted photon use by approximately 15.7% relative to uniform allocation. These results demonstrate the potential of physics-guided sequential sensing to improve photon efficiency in single-photon LiDAR reconnaissance, while further target-absent and experimental validation remain necessary.
